\chapter{Conclusion}
\label{chap:conclusion}

This report documented the research methodology, engineering requirements, system architecture, implementation deliverables, and verification evidence for Project \textbf{FORGE} (\textit{Personalized Procedural Content Generation via Reinforcement Learning}, Project ID: F26-012-FORGE) at the FYP-1 Mid-Term Evaluation milestone.

\section*{Problem and Research Contributions}
Procedural level generation in video games has historically balanced replayability against an abstract aggregate player average, failing to accommodate individual differences in motor ability, combat pacing, and tactical risk tolerance. Concurrently, existing commercial Dynamic Difficulty Adjustment systems rely on brittle, reactive scalar multipliers that trigger jarring rubber-banding and erode feelings of player competence. 

FORGE addresses this gap with a two-stage decoupled generation architecture:
\begin{enumerate}
    \item \textbf{Stage 1:} Synthesizes discrete dungeon room layouts, producing layouts that are connected by construction (random walk with branching), with reachability checked by a breadth-first-search regression test over 3,500 seeded layouts in CI.
    \item \textbf{Stage 2:} Ingests continuous, unprompted in-session telemetry (clearing time, death count, damage ratio) into a Proximal Policy Optimization (PPO) policy to output a continuous 7-dimensional combat difficulty action vector ($\vec{E}, \vec{M}$).
\end{enumerate}
With a pure-GDScript neural network forward pass that runs once per floor in well under 1\,ms, FORGE needs no external runtime and shows that a learned difficulty policy can run deterministically on client hardware. Its closest prior work is experience-driven and player-adaptive PCGRL for platformers~\cite{shu2021edrl, wang2022oparl}; FORGE's contribution is applying closed-loop RL adaptation to several combat and environment attributes inside a shipped roguelike, seeded by an in-game calibration phase.

\section*{Mid-1 Implementation and Verification Achievements}
During Iterations 1 and 2 (Weeks 1 through 10), the project team addressed seven of the committee's ten directives fully and three partially (the remaining acceptance criteria, the empirical judge model, and a demonstrated benefit over a non-adaptive baseline, all scheduled for FYP-1 Final):
\begin{itemize}
    \item Improved visual aesthetics by deploying 2D pixel-art sprites, trauma-based non-linear camera shake, floating combat damage indicators, and professional typography.
    \item Expanded combat variety by introducing Melee Crawlers, kiting Ranged Pyromancers with telegraphed fireball projectiles, and a player dodge roll with invulnerability frames.
    \item Addressed the cold-start problem with a 5-floor calibration sequence that computes an Initial Combat Rating ($CR_0 \in [10, 95]$) with escalating weights ($\vec{w} = [1.0, 1.5, 2.0, 2.5, 3.0]$) and seeds the first PPO observation on the policy's training scales.
    \item Provided a deterministic $+8\%$ per-floor Static Roguelike Baseline mode toggle in the live Defense Telemetry Overlay as the non-adaptive comparison condition.
    \item Built an automated GitHub Actions CI/CD pipeline on Godot 4.7.2 with a headless load check, a 3,500-layout reachability regression test and headless GDScript tests, run on every push to the main branches and every pull request.
\end{itemize}

\section*{Path Forward into FYP-2}
All automated tests pass on the submitted code, and the heuristic fallback keeps the game running without policy weights. A pre-submission audit found and fixed five defects, most importantly unrecorded damage telemetry and a mis-scaled Floor 6 observation. Simulation shows that the trained policy beats an untrained one but personalizes only part of its action space, and frame rate has not yet been logged. The next phases are: correcting the simulator and retraining the policy, then expanding the combat ecosystem to 5 archetypes and training an empirical supervised Judge regressor with TreeSHAP feature attribution at Gate 3 Exit (Final-1); deploying the equipment shop economy, system integration and benchmarking at Gate 4 (Mid-2); and conducting a formal 10-participant human User Acceptance Testing study with GEQ and PENS psychometrics at Gate 5 (Final-2). Project FORGE is on schedule for its Mid-1 milestones, with the open technical risks identified and planned for.
